\documentclass[10pt,a4paper]{amsart}
\usepackage[margin=19mm]{geometry}
\usepackage{amsmath,amssymb,mathtools}
\usepackage[scr=rsfs,cal=euler]{mathalfa}
\usepackage{xfrac}
\usepackage[unicode,hidelinks,bookmarks=false]{hyperref}
\newcommand{\ie}{i.e.\ }
\newcommand{\cf}{cf.\ }
\newcommand{\resp}{respectively\ }
\newcommand{\Subsection}{\subsection}
\providecommand{\nobreakdash}{\nobreak\hbox{-}}
\setlength{\emergencystretch}{2em}
\multlinegap0pt
\usepackage{enumerate}
\usepackage{amssymb}
\usepackage{xcolor}
\newtheorem{lemma}{Lemma}
\newcommand{\D}{\nabla}
\newcommand{\Om}{\Omega}
\newcommand{\R}{\mathbb{R}}
\title{A singular profile for the relativistic heat cost\\ and the special Lagrangian curvature equation}
\author{Xiao-Tian Wu}
\date{Source: arXiv:2607.26970v1 (CC BY 4.0)}
\begin{document}
\maketitle

\noindent\emph{Source and licence.} A singular profile for the relativistic heat cost\\ and the special Lagrangian curvature equation. Version arXiv:2607.26970v1 (CC BY 4.0), distributed by arXiv under the Creative Commons Attribution 4.0 International licence, http://creativecommons.org/licenses/by/4.0/, as recorded in the arXiv licence metadata for this exact version and shown on the abstract page https://arxiv.org/abs/2607.26970v1. Full licence notice: \texttt{licenses/arxiv-2607.26970-CC-BY-4.0.txt}.\par\smallskip

\noindent\textit{Editorial scope.} Counted unit: the article's lemma lem:gs (source0.tex lines 700--702). The article's complete original proof of it is delivered with it (source0.tex lines 703--782, one \texttt{proof} environment of 80 lines). No mathematics is written, repaired or re-derived: the statement and the proof are the article's own lines, in the article's own notation, and no retained line is substituted or re-flowed.  Retained uncounted as the source-local context the proof needs: lines 158--163; lines 387--389; lines 400--405; lines 489--491; lines 493--496; lines 537--545; lines 574--582; lines 633--636; lines 651--666.  Nothing was omitted from any retained span, so the package carries no omission record.  No retained line contains a citation, so the wrapper carries no bibliography and no bibliography entry.  No retained line contains a figure environment, an image-inclusion command or drawing code, so no image is delivered and the package declares no asset dependency.  Editorial formatting only, in addition: an AMS \texttt{amsart} preamble that restates the article's own declarations and macros used by the retained lines, namely the environments \texttt{lemma} on separate counters, together with the standard packages those lines need.  The retained environments print in this document as Lemma 1 in order of appearance, whereas the article prints the same items with the numbering of its own full document.  The complete native source of the article as distributed in the arXiv e-print for this exact version is bundled unchanged as \texttt{sources/206306/source0.tex}.
\begin{equation}\label{eq:ot-pde}
  \det\left(
    \frac{a}{\sqrt{1+|\D u|^2}}
    \left(I-\frac{\D u\otimes\D u}{1+|\D u|^2}\right)D^2u+I
  \right)=\lambda.
\end{equation}

    \begin{equation}\label{eq:T_u-id}
        |T_u(E)|=\lambda|E|\qquad\text{for every Borel set }E\subset\Om.
    \end{equation}

\begin{lemma}\label{lem:in-CS}
    Let $p,q>0$ and $v,w\in\R^{n-1}$ satisfy $p\geq|v|$ and $q\geq|w|$. Then
    \begin{equation}\label{eq:in-CS}
        \sqrt{(p+q)^2-|v+w|^2}\geq\sqrt{p^2-|v|^2}+\sqrt{q^2-|w|^2}.
    \end{equation}
\end{lemma}

\begin{equation}\label{s3:e1}
    T_u(x',t)=(A(t)x',B(t)),
\end{equation}

\begin{equation}\label{eq:A-B}
    A(t)=1-\frac{a}{r\sqrt{1+\dot{r}^2}},\qquad
    B(t)=t+\frac{a\dot{r}}{\sqrt{1+\dot{r}^2}}.
\end{equation}

\begin{equation}\label{eq:ode-sys}
    \left\{
        \begin{aligned}
            \frac{dt}{ds}&=\frac{a\mathcal{A}^{n-1}}{(\lambda-\mathcal{A}^{n-1})(1+s^2)^{3/2}},\\
            \frac{dR}{ds}&=s\frac{a\mathcal{A}^{n-1}}{(\lambda-\mathcal{A}^{n-1})(1+s^2)^{3/2}},\\
            t(0)&=0,R(0)=a.
        \end{aligned}
    \right.
\end{equation}

\begin{lemma}\label{lem:inverse}
    After possibly decreasing $\delta>0$, the function $s\mapsto t(s)$ is strictly increasing on $(-\delta,\delta)$.
    Then there exists a unique inverse function $t\mapsto s(t)$ defined on some interval $(-T,T)$.
    Moreover, the function $s$ satisfies the following asymptotic formula near $t=0$:
    \begin{equation}\label{eq:s-asymp}
        s(t)=\gamma\operatorname{sgn}(t)|t|^{\frac{1}{2n-1}}+o(|t|^{\frac{1}{2n-1}}),
    \end{equation}
    where $\gamma$ is a positive constant depending only on $n,a,\lambda$.
\end{lemma}

\begin{lemma}\label{lem:Ra}
    For every $t\in(-T,T)$ one has $r(t)\geq a$ and $A(t)\in[0,1)$, where $A$ is the radial
    factor in \eqref{eq:A-B}.
\end{lemma}

\begin{lemma}[Choice of radius]\label{lem:radius}
    There exists $\eta\in(0,\min\{T,a/8\})$ such that
    \begin{equation}\label{eq:R-cond}
        \frac{a\,|s(t)|}{\sqrt{1+s(t)^2}}\leq\frac a4\qquad\text{for all }t\in[-\eta,\eta].
    \end{equation}
    Fix such an $\eta$ and write $B_\eta:=B_\eta(0)\subset\R^n$. Then fixing an arbitrary point $(x'_0,t_0)\in B_\eta$, with the notations
    $r_0=r(t_0)$, $A_0=A(t_0)$ and $B_0=B(t_0)$:
    \begin{enumerate}
        \item[\rm(i)] the function $u(x',t):=\sqrt{r(t)^2-|x'|^2}$ is well-defined and positive
        on $B_\eta$;
        \item[\rm(ii)] $|t-B_0|\leq a/2$, and hence $a^2-(t-B_0)^2\geq 3a^2/4$, for all
        $t,t_0\in[-\eta,\eta]$;
        \item[\rm(iii)] $\sqrt{a^2-(t-B_0)^2}\geq|x'-A_0x_0'|$ for all
        $(x',t),(x_0',t_0)\in B_\eta$.
    \end{enumerate}
\end{lemma}

\begin{lemma}\label{lem:gs}
    The function $u$ is a generalized solution to \eqref{eq:ot-pde}.
\end{lemma}

\begin{proof}
    {\it Step 1: $u$ is $c$-convex in $B_\eta$.}

    It suffices to show that for an arbitrary point $(x_0',t_0)\in B_\eta$,
    there exists a $c$-support of $u$ at $(x_0',t_0)$.
    Note that $T_u(x',t)=(A(t)x',B(t))$, where $A$ and $B$ are given in \eqref{eq:A-B}.
    For simplicity of notation, set $r_0=r(t_0)$, $A_0=A(t_0)$ and $B_0=B(t_0)$.
    We will establish a one-dimensional support inequality for the function $r(t)$ and then obtain the support inequality for $u(x',t)$ by virtue of Lemma~\ref{lem:in-CS}.

    By Lemma~\ref{lem:radius}(ii), the function $g(t):=A_0r_0+\sqrt{a^2-(t-B_0)^2}$ is well-defined on $(-\eta,\eta)$.
    Define $h(t):=r(t)-g(t)$.
    Then $h(t_0)=\dot{h}(t_0)=0$.
    We see that $\dot{g}$ is strictly decreasing because $\ddot{g}(t)=-a^2/(a^2-(t-B_0)^2)^{3/2}<0$.
    On the other hand, $\dot{r}=s$ is strictly increasing by Lemma~\ref{lem:inverse}.
    Putting these facts together, we further see that $h$ attains its minimum at $t=t_0$.
    It follows that for every $t\in(-\eta,\eta)$ there holds
    \begin{equation}\label{eq:r-support}
        r(t)\geq A_0r_0+\sqrt{a^2-(t-B_0)^2}.
    \end{equation}

    Since $A_0\geq0$ by Lemma~\ref{lem:Ra} and $r_0\geq|x_0'|$, we have $A_0r_0\geq A_0|x_0'|$.
    On the other hand, by Lemma~\ref{lem:radius}(iii), $\sqrt{a^2-(t-B_0)^2}\geq|x'-A_0x_0'|$ on $B_\eta$.
    Then
    \begin{align*}
       &u(x',t)-\Big(u(x_0',t_0)+c((x',t),T_u(x_0',t_0))-c((x_0',t_0),T_u(x_0',t_0))\Big)\\
       =&\sqrt{r(t)^2-|x'|^2}-A_0\sqrt{r_0^2-|x_0'|^2}-\sqrt{a^2-|x'-A_0x_0'|^2-(t-B_0)^2}\\
       \geq&\sqrt{\Big(A_0r_0+\sqrt{a^2-(t-B_0)^2}\Big)^2-|A_0x_0'+(x'-A_0x_0')|^2}\\
       &-\sqrt{(A_0r_0)^2-|A_0x_0'|^2}-\sqrt{a^2-(t-B_0)^2-|x'-A_0x_0'|^2}\geq0,
    \end{align*}
    where \eqref{eq:r-support} is used in the second-to-last inequality and Lemma~\ref{lem:in-CS} is used in the last inequality.
    Define
    \[
    \ell(x',t):=u(x_0',t_0)+c((x',t),T_u(x_0',t_0))-c((x_0',t_0),T_u(x_0',t_0)).
    \]
    Then $\ell(x_0',t_0)=u(x_0',t_0)$ and $\ell\leq u$ on $B_\eta$, namely $\ell$ is a $c$-support of $u$ at $(x_0',t_0)$.
    Since $(x_0',t_0)$ is arbitrary, we conclude that $u$ is $c$-convex in $B_\eta$.

    {\it Step 2: $u$ satisfies \eqref{eq:T_u-id} for every Borel set $E\subset B_\eta$.}

    Fix an arbitrary Borel set $E\subset B_\eta$.
    Recall that $T_u$ is a single-valued map given by \eqref{s3:e1}.
    We reparametrize $T_u$ via $s$.
    For simplicity of notation, set
    \begin{align*}
        &\mathcal{A}(s):=A(t(s))=1-\frac{a}{R(s)\sqrt{1+s^2}},\\
    &\mathcal{B}(s):=B(t(s))=t(s)+\frac{as}{\sqrt{1+s^2}}.
    \end{align*}
    Then $T_u(x',t(s))=(\mathcal{A}(s)x',\mathcal{B}(s))$.
    Since $t'\geq 0$ and $a>0$, we have
    \[
    \mathcal{B}'(s)=t'+\frac{a}{(1+s^2)^{3/2}}>0.
    \]
    Hence $\mathcal{B}$ is a strictly increasing smooth function on $(s(-\eta),s(\eta))$.
    Let $\mathcal{B}^{-1}$ denote the inverse function of $\mathcal{B}$.

    We calculate $|T_u(E)|$ by slicing.
    Given a height $w\in(\mathcal{B}(s(-\eta)),\mathcal{B}(s(\eta)))$, the $w$-slice of $T_u(E)$ is given by
    \[
    (T_u(E))_w:=\{x'\in\R^{n-1}:(x',w)\in T_u(E)\}=\mathcal{A}(\mathcal{B}^{-1}(w))E_{t(\mathcal{B}^{-1}(w))},
    \]
    where $E_{t(\mathcal{B}^{-1}(w))}:=\{x'\in\R^{n-1}:(x',t(\mathcal{B}^{-1}(w)))\in E\}$.
    It follows that
    \begin{equation}\label{eq:slice}
        |(T_u(E))_w|=(\mathcal{A}(\mathcal{B}^{-1}(w)))^{n-1}|E_{t(\mathcal{B}^{-1}(w))}|.
    \end{equation}
    Using the first and the second equations in \eqref{eq:ode-sys}, we obtain
   \begin{equation}\label{eq:ode-id}
    (\mathcal{A}(s))^{n-1}\mathcal{B}'(s)=\lambda t'(s),\quad\forall s\in(-s(\eta),s(\eta)).
   \end{equation}
   Using \eqref{eq:slice} and \eqref{eq:ode-id} in order, we have
   \begin{align*}
    |T_u(E)|=&\int_{\mathcal{B}(s(-\eta))}^{\mathcal{B}(s(\eta))}|(T_u(E))_w|dw\\
    =&\int_{\mathcal{B}(s(-\eta))}^{\mathcal{B}(s(\eta))}(\mathcal{A}(\mathcal{B}^{-1}(w)))^{n-1}|E_{t(\mathcal{B}^{-1}(w))}|dw\\
    =&\int_{s(-\eta)}^{s(\eta)}(\mathcal{A}(s))^{n-1}|E_{t(s)}|\mathcal{B}'(s)ds\\
    =&\lambda\int_{s(-\eta)}^{s(\eta)}|E_{t(s)}|t'(s)ds=\lambda|E|.
   \end{align*}
   Since $E$ is arbitrary, we conclude that \eqref{eq:T_u-id} holds for every Borel set in $B_\eta$.

   Putting Step 1 and Step 2 together, Lemma~\ref{lem:gs} is proved.
\end{proof}

\end{document}
